% HAND-WRITTEN fragment -- NOT generated by maestro2tex.
% Source: ~/vestarv/docs/afe_rev2_pixel_verify.html rev. 2026-08-23, S7.3 and
% S8.2 (run simruns/pixtop2_gain, dc sweep ncode 0..319 step 1, inert electrode
% cox = cred = 0, nominal, 37 C):
%   V(WE) - V_CM = -619.6424 uV at N = 0 and -619.6421 uV at N = 319, i.e.
%     0.3 nV of spread over all 320 codes;
%   V_OUT - V_CM = -619.6444 uV (N = 0) ... -623.5525 uV (N = 319), a 3.908 uV
%     drift over dRf = 1.914 MOhm -> 2.042 pA of A2 input current;
%   i_off = -34.4247 nA (N = 0), -1.5651 nA (N = 63), -0.3227 nA (N = 319).
% The 2 nA input-referred noise budget is the in_rma spec used throughout this
% tree (see ss:tiarprog).

\subsubsection{Zero-Current Axis}\label{sss:pixtop-zero}

With the electrode inert the working electrode sits \SI{619.64}{\micro\volt}
below the common mode, and it is the same \SI{619.64}{\micro\volt} at every one
of the 320 gain codes: \SI{0.3}{\nano\volt} of spread. No current flows in
\(R_f\), so the output carries A2's input offset alone and nothing that scales
with the tap. \textbf{One zero measurement therefore calibrates every gain
code} --- the result the per-gain-code current zero of
Table~\ref{tab:calibration-constants} rests on, and the reason the common-mode
slot on the converter multiplexer is worth its silicon.

What survives that calibration is input bias current. The output drifts a
further \SI{3.91}{\micro\volt} between \(N = 0\) and \(N = 319\) while \(R_f\)
grows by \SI{1.914}{\mega\ohm}: \SI{2.04}{\pico\ampere} into A2's inverting
input, \SI{0.5}{\percent} of the \SI{0.41}{\micro\ampere} full scale at maximum
gain and negligible at every code below it.

Referred to current, the uncalibrated zero is \SI{34.4}{\nano\ampere} at
\(N = 0\), \SI{1.57}{\nano\ampere} at \(N = 63\) and
\SI{0.323}{\nano\ampere} at \(N = 319\). The last two are already inside the
\SI{2}{\nano\ampere} input-referred noise budget, so at high gain the
uncalibrated offset is not what limits a reading; at the bottom code it is
\SI{34.4}{\nano\ampere} on a \SI{40.7}{\micro\ampere} full scale and below one
converter code either way (Section~\ref{sss:pixtop-adc}).
